\documentclass[11pt]{amsart}
\usepackage[margin=1.1in]{geometry}
\usepackage{amsmath,amssymb,amsthm,mathtools}
\usepackage{lmodern}
\usepackage[mathscr]{eucal}
\usepackage{booktabs,enumitem,xcolor,array}
\usepackage[colorlinks=true,linkcolor=blue!50!black,citecolor=green!40!black,urlcolor=blue!60!black]{hyperref}
\hypersetup{
  pdftitle={Positive solutions of the explicit formula for the Riemann zeta function: near-criticality and uniqueness},
  pdfauthor={Nic Johns},
  pdfkeywords={Riemann zeta function, Weil explicit formula, positivity, linear programming duality, conductor bounds, Fourier interpolation, Hermite interpolation, computer-assisted proof}
}

\setlength{\emergencystretch}{2.5em}
% macros.tex -- shared macros for main.tex and sections/*.tex.
% Notation is introduced in the text (Sections 1.6, 2.1, 3 and 4); do not introduce competing symbols.
% New local macros in a section file: use \providecommand and a section-specific prefix.

% ---------------------------------------------------------------- general
\newcommand{\RR}{\mathbb R}
\newcommand{\ZZ}{\mathbb Z}
\newcommand{\CC}{\mathbb C}
\newcommand{\QQ}{\mathbb Q}
\newcommand{\dd}{\,\mathrm d}
\newcommand{\eps}{\varepsilon}
\renewcommand{\Re}{\operatorname{Re}}
\renewcommand{\Im}{\operatorname{Im}}
\newcommand{\supp}{\operatorname{supp}}
\newcommand{\one}{\mathbf 1}
\newcommand{\FT}[1]{\widehat{#1}}          % Fourier transform, \int f(t)e^{-2\pi i\xi t}dt
\newcommand{\hF}{\widehat F}
\newcommand{\hyp}[1]{\textup{(#1)}}       % named hypothesis tag, e.g. \hyp{N}, \hyp{Mg}
\newcommand{\astranote}{\footnote{Due to Astra (OpenAI), contributed during an independent review of an earlier version of this
  paper, and independently verified.}}

% ---------------------------------------------------------------- coordinates
\newcommand{\xin}[1]{\xi_{#1}}            % \xi_n = \log n/(2\pi)

% ---------------------------------------------------------------- arithmetic data
\newcommand{\PP}{\mathrm{PP}}             % prime powers > 1; \PP_J the first J
\newcommand{\xiPP}{\xi_{\mathrm{PP}}}     % {\xi_n : n \in PP}
\newcommand{\BRS}{\mathcal M}             % {\log m/(4\pi) : m \ge 1}
\newcommand{\GR}{\Gamma_{\mathbb R}}
\newcommand{\GC}{\Gamma_{\mathbb C}}
\newcommand{\Xiinf}{\Xi_\infty}           % archimedean part of \Xi
\newcommand{\Zz}{Z_\zeta}                 % real zero ordinates of zeta
\newcommand{\muz}{\mu_\zeta}
\newcommand{\nuz}{\nu_\zeta}
\newcommand{\pz}{p_\zeta}

% ---------------------------------------------------------------- framework
\newcommand{\Oinf}{\Omega_\infty}         % archimedean density
\newcommand{\Arch}{\mathcal A}            % archimedean functional; \Arch_q with conductor q
\newcommand{\Tc}{\mathcal T}              % test class; \Tc_\delta
\newcommand{\Gc}{\mathcal G}              % Gaussian wave packets
\newcommand{\Spec}{\mathcal R}            % right side of the explicit formula, \Spec_{\mu,\nu}(F)
\newcommand{\Kad}{\mathcal K}             % admissible pairs
\newcommand{\Cone}{\mathcal C}
\newcommand{\ConeOPS}{\mathcal C_{\mathrm{OPS}}}
\newcommand{\kstar}{\kappa^{*}}
\newcommand{\kOPS}{\kappa^{*}_{\mathrm{OPS}}}
\newcommand{\qmin}{q_{\min}}
\newcommand{\condE}{\textup{(E)}}
\newcommand{\condU}{\textup{(U)}}
\newcommand{\condS}{\textup{(S)}}
\newcommand{\RH}{\mathrm{RH}}

% ------------------------------------------------------------- zero-killing family
\newcommand{\Wnull}{W_0}                  % z^2[(z^2+9)K_0 - 6zK_1]
\newcommand{\Psione}{\Psi_1}              % Gamma-only kernel
\newcommand{\Hc}{\mathscr H}              % growth class \Hc_a
\newcommand{\Jset}{\mathcal J}            % {J : M_J non-singular}
\newcommand{\Frep}{F_{\mathrm{rep}}}      % certified rational representative



% ---------------------------------------------------------------- environments
% theorem/proposition/lemma/corollary: proved in the paper (possibly computer-assisted,
% with the computation and ancillary file named).  conjecture: open.
% observation ("Numerical observation"): computed, not certified, never used in a proof.
\newtheorem{theorem}{Theorem}[section]
\newtheorem{proposition}[theorem]{Proposition}
\newtheorem{lemma}[theorem]{Lemma}
\newtheorem{corollary}[theorem]{Corollary}
\newtheorem{conjecture}[theorem]{Conjecture}
\newtheorem{question}[theorem]{Question}
% Lettered conjectures: \begin{lconj}{U}\label{conj:U} ... \end{lconj} prints "Conjecture U",
% and \ref{conj:U} prints "U".  Used for Conjecture U and Conjecture S only.
\newtheorem{lconjinner}{Conjecture}
\newenvironment{lconj}[1]{\renewcommand{\thelconjinner}{#1}\begin{lconjinner}}{\end{lconjinner}}
% Lettered headline theorems: \begin{lthm}{A} ... \end{lthm} prints "Theorem A".
\newtheorem{lthminner}{Theorem}
\newenvironment{lthm}[1]{\renewcommand{\thelthminner}{#1}\begin{lthminner}}{\end{lthminner}}
% Unnumbered corollary for the introduction.
\newtheorem*{introcor}{Corollary}
\theoremstyle{definition}
\newtheorem{definition}[theorem]{Definition}
\theoremstyle{remark}
\newtheorem{remark}[theorem]{Remark}
\newtheorem{observation}[theorem]{Numerical observation}
\numberwithin{equation}{section}

% Section files; paths are relative to the top-level directory (arXiv compiles from there).
\newcommand{\inputsection}[1]{\input{sections/#1}}

\title[Positive solutions of the explicit formula]{Positive solutions of the explicit formula for $\zeta(s)$:\\ near-criticality and uniqueness}

\author{Nic Johns}

\date{October 2026}

\subjclass[2020]{Primary 11M26; Secondary 11M06, 11R29, 11R42, 41A05, 42A82, 65G20, 90C05}
\keywords{Riemann zeta function, Weil explicit formula, positivity, linear programming duality, conductor bounds, Fourier interpolation, Hermite interpolation, computer-assisted proof}

\begin{document}

\begin{abstract}
We study the positive solutions of Weil's explicit formula with the archimedean data of the Riemann zeta function:
pairs of positive measures, one on the zero side and one on the prime side beyond $\log2/2\pi$. By linear-programming
duality such pairs exist if and only if a positivity statement involving neither zeros nor primes holds; its slack
$\kappa^*$ gives the optimal degree-one explicit-formula conductor bound $q_{\min}=e^{-2\pi\kappa^*}$. We conjecture that $\zeta$'s
own pair is the only solution; this does not imply the Riemann hypothesis, but makes it equivalent to existence. We
prove the conjecture for pairs carried by the zeros of $\zeta$ plus finitely many points and by the
Bondarenko--Radchenko--Seip nodes, using positivity of the prime measure only on one residue class of nodes, and of the zero measure only at the origin. Certificates of the form
$\Xi^2(H+\varepsilon e^{-\pi t^2})$, with $H$ a polynomial in $t^2$, vanish at every zero of $\zeta$ and show
unconditionally that $\kappa^*\le1.1508391\cdot10^{-906}$. Uniqueness would follow from asymptotic properties of a
family of such functions, defined by Hermite interpolation at the prime powers.
\end{abstract}

\maketitle

\setcounter{tocdepth}{1}
\tableofcontents

\inputsection{01-intro}
\inputsection{02-setting}
\inputsection{03-uniqueness}
\inputsection{04-family}
\inputsection{05-certified}
\inputsection{06-conjectures}

\section*{AI-use statement}
This paper was produced with substantial use of AI. The mathematical arguments, the numerical computations and the text were
generated by instances of Claude Opus 5.5 (Anthropic) working as coordinated agents: one set of agents developed the arguments
(building on \cite{BRS} and on the explicit-formula method for discriminant and conductor bounds \cite{Sta74,Odl76,Ser75,Poi77,Poi77b}), separate agents
acted as independent hostile referees of each component, and the referees' requested repairs were incorporated. The numerical
certificates were computed by AI-written scripts that are included as supplementary material (Appendix~\ref{app:anc}) and can be
rerun. The author initiated and directed this project, designed its verification process, and reviewed the results and the
verification record, but did not independently check every step of the written proofs. The author takes full responsibility for the
content of this paper, including any errors, and welcomes independent verification and corrections. Section~\ref{ssec:claims}
states precisely what is claimed.

\appendix
\inputsection{A-certification}
\inputsection{B-proofs}
\inputsection{C-anc}

\bibliographystyle{production/amsplain-doi}
\bibliography{refs}

\end{document}
